\hypertarget{an-operational-ontology-for-four-banking-journeys}{%
\section{An operational ontology for four banking
journeys}\label{an-operational-ontology-for-four-banking-journeys}}

An ontology is useful when it makes different interpretations explicit
before they enter an estimate. The proposed causal ontology should
distinguish an entity from an observation about that entity, a concept
from a measured variable, a variable definition from a particular
observation, and an intervention from the administrative event used to
record it. These distinctions are not cosmetic. They determine which
quantities can be assigned, which can be conditioned on, which share
measurement errors and which belong to the same population.

For the funding journey, consider ``interest rate.'' A contractual
annual rate, a monthly interest payment, a competitor's advertised rate
and a fitted marginal funding cost are different variables. A unit
conversion can relate two encodings of the same rate, but it cannot turn
a payment into an offer without account balance, accrual and timing
information. Nor does it turn an advertised offer into a received offer.
The ontology should identify the product, currency, contractual basis,
decision date and exposure window. An agent can use a glossary to
suggest the mapping; the analyst must still see which distinctions the
mapping retains.

The associated entity hierarchy also matters. An account belongs to a
legal customer, a customer may belong to a household, and several
records may refer to the same account. Those relations do not themselves
say that one node causally influences another. They do affect
aggregation and dependence. An offer randomized by household cannot be
analyzed as though each account were assigned independently merely
because the extract has one row per account. If a latent household
membership is replaced by a record-linking prediction, that substitution
introduces the same kind of measurement question developed in the proxy
example.

For sequential servicing, ``treatment'' may mean eligibility,
assignment, receipt, completion or adherence. The stages can occur on
different dates and have different causes. A rule that assigns
assistance after an arrears flag targets a dynamic regime; comparing
customers who eventually complete assistance targets a selected group
unless further assumptions justify a different interpretation. The
ontology should retain the action that the bank can choose and the later
events that it can only influence. Otherwise an assistant may propose an
infeasible intervention on a customer's future response while describing
it as an ordinary policy choice.

For discovery, a variable needs a measurement protocol as well as a
label. A repeated risk score can change because the customer changed,
because the scoring model was updated or because new source information
became available. Treating all vintages as the same mechanism can make
an apparent causal change coincide with a measurement change. A diagram
can represent the score-generating mechanism or restrict the analysis to
a stable declared regime. The correct choice depends on the question; a
vocabulary match alone cannot choose it.

For trustworthy observations, separate an event from its record, a
proposed match from a true identity, and label correctness from the
value being labeled. A classifier's probability that a record pair
matches is an algorithm output with a calibration question. It is not an
observed identity. A hand-reviewed match can supply additional
information, but its reviewer, selection process and uncertainty remain
relevant. The ontology can preserve these distinctions in a public
synthetic example long before the proposal needs an expensive bank
labeling exercise.

Across all four journeys, the proposed variable definition therefore
contains a conceptual role, unit of analysis, value domain, measurement
mechanism, clocks, lineage and relation to a feasible intervention. The
system should support unresolved fields explicitly. A missing clock
should prompt a targeted evidence request; an unresolved proxy relation
should produce alternative measurement models. Automatically completing
every field with plausible text would hide the very uncertainty that a
causal proposal must expose.

\hypertarget{semantic-relations-need-typed-boundaries}{%
\section{Semantic relations need typed
boundaries}\label{semantic-relations-need-typed-boundaries}}

Common ontology relations carry different logical meanings. ``A savings
account is an account'' describes a taxonomy. ``A household owns an
account'' describes a relation between entities. ``A risk score is
computed from balance history'' describes a derivation. ``An offer
changes a future balance'' proposes a causal mechanism. A general
relation extractor may return all four as triples, but a causal engine
cannot assign them all the same edge semantics.

A taxonomy helps establish which definitions apply to an instance; it
does not imply that intervening on an instance's category has a coherent
economic meaning. Ownership determines which exposures should be
combined and may enter a separate causal model, but the ownership triple
alone does not identify an intervention effect. A derived score carries
deterministic or stochastic dependence on its inputs; treating the score
and every input as unrelated primitive variables can contradict the
declared derivation. Each relation needs a type before a diagram
operation acts on it.

Duplicate encodings provide a simple formal example. Suppose a decimal
rate \(R\) and its basis-point encoding \(B\) satisfy \(B=10{,}000R\)
exactly. A vocabulary matcher may list them as two numerical variables.
Their observational distribution lies on a deterministic line. An
imagined intervention setting \(R=r\) and \(B=b\) independently is
incompatible with the encoding relation unless \(b=10{,}000r\), or
unless the target is explicitly a system fault that breaks that
relation. Ordinary policy analysis should treat these as two coordinates
of one declared quantity, not as two independently selectable business
levers.

Alias resolution must be reversible. Two fields with the same display
name can describe different vintages; two differently named fields can
be identical encodings. The system should retain source identifiers,
proposed canonical identifiers and the transformation between them. A
bijective unit conversion can be checked mechanically. A many-to-one
coarsening, such as mapping several delinquency states into ``at risk,''
discards information and can change the adjustment or decision problem.
It therefore needs a scientific interpretation as well as a data
transformation.

The preferred local baseline is a deterministic dictionary and lineage
validator. It can detect an undeclared conversion, missing unit,
duplicate identifier or a cycle in a purported derivation. A language
assistant adds value if it finds plausible semantic conflicts that the
dictionary did not encode and supplies reviewable evidence for resolving
them. The two should be compared on missed consequential distinctions
and unnecessary review work. Merely producing more complete metadata is
not enough when some of the added metadata is wrong.

The Selby study provides relevant adverse evidence for automatic
completion. Its inspected experiments elicited quantitative priors and
imputed tabular values from language models. Repeated prompts often
produced consistent responses, yet some priors conflicted strongly with
observations and conventional nearest-neighbor or random-forest
imputation often performed better. The authors identify ambiguous units,
incomplete descriptions and recognizable public datasets as limitations.
This supports careful metadata and benchmark design. It does not
establish a theorem about causal identification or license silent
imputation of a bank's missing assignment mechanism.
\citep{v12_la_2402_07770}

\hypertarget{a-treatment-label-can-hide-several-interventions}{%
\section{A treatment label can hide several
interventions}\label{a-treatment-label-can-hide-several-interventions}}

Consider an assumed assistance program with two versions. Without
assistance, successful repayment has probability .4. The low-intensity
version has success probability .42 and costs USD 10; the high-intensity
version has success probability .6 and costs USD 100. Suppose assignment
and version selection are randomized in this constructed example, and
each additional successful repayment is assigned a planning value of USD
1,000. The version probabilities define the program being evaluated.

If \(w\) is the probability of the high-intensity version among treated
accounts, the effect of that mixed program is

\[
 \tau(w)=(1-w)(.42)+w(.6)-.4=.02+.18w.
\]

Its expected cost is \(c(w)=10+90w\), so the assumed net value is

\[
 v(w)=1000\tau(w)-c(w)=10+90w.
\]

At \(w=1/4\), the effect is .065 and the net value is USD 32.50. At
\(w=3/4\), the effect is .155 and the net value is USD 77.50. Every
calculation is exactly reproduced from the stipulated rational
probabilities. No observed banking profitability is asserted.
\citep{exec13_ontology_examples}

Both datasets could contain the same binary column named ``assistance.''
Nevertheless, the programs have different meanings and different costs.
An assistant should ask whether the deployment preserves the version mix
and the mechanisms under which those version-specific outcomes were
obtained. Reporting a precisely estimated effect for the earlier mix as
the effect of the later mix would answer the wrong question even if the
graph's two main node labels were unchanged.

The construction is deliberately simple. If version choice depends on
baseline need, a dynamic or covariate-dependent mixture \(w(x)\) is
relevant. If customers select their own version after assignment,
assignment and receipt become different interventions. If staff
implement nominally identical assistance differently, the declared
version may still omit consequential variation. A useful ontology does
not promise to remove all such ambiguity; it identifies which variation
is included in the intervention definition and which is left as an
unresolved transport assumption.

This also explains why a prototype should retain version information
before it is needed by an estimator. A source package might accept one
binary treatment and correctly estimate the target supplied to it. The
proposal's common layer is responsible for making that supplied target
faithful to the actual decision. An LLM can identify the hidden-version
possibility and help locate policy documentation. The effect arithmetic
then belongs to a separately specified calculation, with the assumed
assignment and valuation law visible to the reader.

\hypertarget{a-restricted-ontolearner-integration}{%
\section{A restricted OntoLearner
integration}\label{a-restricted-ontolearner-integration}}

OntoLearner supplies an available component for the semantic stage. Its
paper describes ontology loading, learning tasks and learner models,
including typing, taxonomy and relation discovery. It also reports
domain-dependent error patterns, retrieval limitations and calibration
differences. The inspected requirements were defined by the developers,
and counts of ontologies or requirements differ across passages. A
proposal can use the architecture and reported limitations without
turning those counts into a certified inventory or treating every
extracted relation as causal. \citep{v12_la_2607_01977}

The inspected repository revision is
074a41ba23a85f5b72d4e28da042eab3d50bd90a, dated 3 September 2026. Its
loader chooses a local source when a path is supplied and otherwise
calls a Hugging Face download. The inspected download call supplies a
repository, file and dataset type, but no revision. Further extraction
calls download typing, taxonomy and non-taxonomic-relation files without
a revision argument as well. Pinning this Python source therefore does
not by itself pin the ontology bytes that an automatic download will
return. \citep{v12_la_package_ontolearner}

There is another relevant boundary in the imported ontology closure. The
base class's \texttt{contains\_imports} hook returns false. When a
subclass enables that hook, the loader follows resolved imports, but it
catches import-loading failures, logs them and continues; an unresolved
import also produces a warning without failing the load. A nonempty
top-level graph therefore does not establish that all intended imported
definitions were loaded. If an unavailable import contains the
definition distinguishing a legal customer from a household, the
resulting partial ontology could still look operationally successful
while supporting an incomplete semantic map. This is a source-inspected
failure possibility, not a report of an executed OntoLearner trial.
\citep{exec13_ontolearner_inspection}

The proposed integration should admit a fixed local ontology bundle with
explicit revision, byte identities and an import manifest. A wrapper
should report complete, intentionally partial or failed loading,
preserving which definitions are absent. It should carry the original
relation types into a semantic proposal record and require a separate
reviewed mapping before creating causal nodes or edges. The initial
useful output is a candidate variable dictionary with evidence and
conflicts, not an automatically accepted causal graph.

The package's MIT license permits reuse subject to its notice and
disclaimer conditions. Its inspected project file declares Python at
least 3.10 and below 3.14, together with RDF, graph, dataframe,
validation and model-related dependencies; several version
specifications are broad. Reusing a small local loading/extraction path
should still be assessed against a locked actual environment and its
transitive dependencies. A source inspection provides no build, memory
or compatibility measurement. The proposed design is to evaluate a restricted, reproducible semantic path
and defer automatic model-backed ontology expansion until its task, data
rights, resource use and evidence handling have been assessed. The preceding
ontology chapter also shows why document identity must be checked separately
from lexical similarity before its extraction metric can support a causal
observation contract.
\citep{v12_la_package_ontolearner}

Porting the whole library would preserve neither ontology identity nor
causal meaning automatically. A port could be justified later by a
measured platform constraint, but it would require equivalent parsing,
import handling and relation semantics. Direct adoption of a mutable
remote default would make a reproduction depend on an unspecified
ontology version. Rejecting the entire library would discard useful
semantic infrastructure. The restricted path retains its useful
operations while making their scientific handoff explicit.

\hypertarget{evidence-extraction-as-a-separate-experiment}{%
\section{Evidence extraction as a separate
experiment}\label{evidence-extraction-as-a-separate-experiment}}

A claim about an edge has at least three evidential questions. Was the
relevant source obtained and correctly parsed? Does an inspected passage
actually support the narrow proposition attributed to it? Is that
proposition applicable and adequate for the intended banking model? A
citation attached to a graph is useful only if the workflow can answer
these questions separately. Text extraction accuracy cannot establish
causal validity, and a valid causal result from another population does
not automatically establish transport.

The ARGUS study provides a concrete adjacent comparison. It audits
reported evidence for identifying assumptions in
difference-in-differences studies and explicitly does not adjudicate
causal truth. Its flaw-injection experiments supply controlled local
references, while its small human-label pilot diagnoses failures in
real-paper auditing. The gated pipeline's omissions include cases in
which retrieval found no relevant passage and the system abstained. In
the pilot, human reviewers often found evidence for cells where the
system had abstained; the paper's subsequent calibration changes were
evaluated on the same papers. These details matter more to our design
than a single aggregate agreement number. \citep{v12_la_2609_30867}

For a diagram assistant, a suitable first experiment can be public and
inexpensive. Assemble passages whose scope is explicit, paired with
propositions that they support, contradict or leave unresolved. Include
genuine parse failures, omitted context, differing populations and
statements that concern association rather than intervention. An
independent reference should record the exact passage and the reason for
its label. Generated alterations can expose specific failure mechanisms,
but their artificial nature should remain visible and should not be
counted as representative bank documents.

The baseline should include direct passage lookup and a deterministic
field/rule check before introducing a model. A model-based system should
report the retrieved material and the proposition it inferred from it.
If the reference passage was never retrieved, the retrieval component
failed. If it was retrieved but misread, the interpretation failed. If
it was read accurately but promoted beyond its population or
assumptions, the scientific mapping failed. Different remedies follow:
improve acquisition/search, improve task interpretation, or tighten the
rules for accepting a causal premise.

An observed negative result also needs precise wording. ``No relevant
evidence was found in this bounded search'' records an acquisition
outcome. ``The inspected document does not state how assignment
occurred'' records a document-level absence under the inspected scope.
``Assignment was confounded'' is a claim about the underlying process.
The first two do not prove the third. Conversely, a polished statement
that assignment was ``as if random'' does not establish the required
condition merely because the text contains the expected vocabulary.

This experiment can support a decision about building a literature and
documentation assistant before paying for a bank audit. It cannot show
that the bank's declared assignment process is true. The eventual
external input is a process-specific record or experiment, sought only
if the prototype and the investment calculation establish its value. The
staged evidence argument therefore has a concrete boundary: public
documents test extraction and reasoning obligations; target-specific
evidence tests propositions about the target institution.

\hypertarget{reviewing-conflicts-without-counting-votes-as-evidence}{%
\section{Reviewing conflicts without counting votes as
evidence}\label{reviewing-conflicts-without-counting-votes-as-evidence}}

Consider three separately proposed restrictions: an offer changes
engagement, engagement changes a later score, and the score determines
the earlier offer. They cannot all be direct temporal assertions about
the same variable instances in an acyclic model. A syntax checker can
identify the contradiction. It cannot determine whether the score was
actually measured earlier, whether ``changes'' meant association,
whether a date was omitted, or whether a cyclic equilibrium model was
intended. The review should return these alternatives to the analyst
with their evidence.

One useful formal object is the set of admitted models after each
restriction. If \(K_1,K_2,K_3\) are proposed predicates, the combined
class is

\[
 \{m\in\mathcal M:K_1(m)\land K_2(m)\land K_3(m)\}.
\]

If it is empty, the system can search for a smallest subset of the
predicates that is already inconsistent within the declared model class.
Reporting that subset localizes the conflict. It does not rank the
predicates by truth or show that the smallest deletion is scientifically
appropriate. Several minimal conflicting subsets may exist; suppressing
all but the first can leave the analyst unaware of another
contradiction.

Priority should follow the evidence and the decision. A verified
timestamp can refute a purported same-instance temporal order. An
expert's general mechanism description may instead need to be
instantiated at a different date. A model suggestion with no supporting
source may remain a hypothesis. If withdrawing one disputed restriction
changes a robustly positive value into a set containing a substantial
loss, that restriction merits attention before a remote relation whose
alternatives leave the decision unchanged. This is a proposed review
policy, not a claim that the system can mechanically rank every kind of
scientific evidence.

Conflicts between sources also need not be contradictions. Two studies
can investigate different treatment versions, observation regimes or
populations. The ontology and intervention-version records make those
differences explicit. A source that supports an association at branch
level may have little to say about an account-level intervention. A
diagram can preserve both claims at their proper levels rather than
force one to win a debate about an underspecified phrase.

The proposal's deliverable is thus a reviewable sequence from language
to typed variables, explicit premises, compatible models, target
calculations and decision consequences. Each transition has examples
that can pass, fail or remain unresolved. That sequence gives an agent
useful work throughout the causal analysis while preserving the
mathematical and evidential limits of each contribution.
\label{exec13:ontology-development-end}
